% !TeX root = ../complete.tex
\section{مدل تأخیر \lr{RC} و \lr{Elmore Delay}}
\label{sec:delay-rc-model}

برای تخمین دستی تأخیر، ترانزیستور روشن را با یک مقاومت و ظرفیت‌های
وابسته به ولتاژ را با مقادیر نماینده در بازه گذار مدل می‌کنیم.
هدف این مدل، نشان‌دادن اثر مسیر هدایت و بار روی زمان کلیدزنی است.
دقت تخمین به گذار، مقدار مقاومت و ظرفیت‌های انتخاب‌شده وابسته است
\cite[بخش \secnum{۴.۳}، صفحات ۱۴۶ تا ۱۵۵]{weste}.

\subsection{مقاومت، ظرفیت و انتخاب مسیر هدایت}
\label{sec:delay-rc-sizing}

در مدل کتاب، ترانزیستور خاموش مسیر جریان را قطع می‌کند. ترانزیستور
روشن مانند کلیدی بسته با مقاومت مؤثر در مسیر هدایت عمل می‌کند.
اگر عرض افزاره \(k\) برابر عرض واحد باشد و طول کانال و شرایط مدل
ثابت بمانند، تقریب می‌زنیم که جریان \(k\) برابر و مقاومت \(1/k\)
برابر شود. مقاومت واحد \lr{nMOS} را \(R\) می‌نامیم. در قرارداد
آموزشی کتاب، مقاومت واحد \lr{pMOS} برابر \(2R\) است. بنابراین
مقاومت مدل برای \lr{nMOS} و \lr{pMOS} با ضریب عرض \(k\)
به‌ترتیب \(R/k\) و \(2R/k\) است. ضریب ۲ فرض مثال است و به فناوری
وابسته است \cite[بخش \secnum{۴.۳.۱}، صفحه ۱۴۶]{weste}.

اگر \(C\) ظرفیت ورودی ترانزیستور واحد باشد، ظرفیت ورودی ترانزیستور
با ضریب عرض \(k\) را \(kC\) می‌گیریم. هر ناحیه تماس‌دار
\lr{Diffusion} در افزاره واحد نیز در تقریب آموزشی ظرفیتی حدود
\(C\) دارد. ظرفیت‌های روی یک گره جمع می‌شوند. جمله‌های ظرفیت، پس
از انتخاب مسیر شارژ یا تخلیه، باید از روی اتصال‌های همان مدار شمرده
شوند \cite[بخش‌های \secnum{۴.۳.۲} و \secnum{۴.۳.۳}، صفحه ۱۴۷]{weste}.

برای \lr{Inverter} مرجع با عرض‌های نسبی ۱ برای \lr{nMOS} و ۲
برای \lr{pMOS}، ورودی مشترک ظرفیت \(C+2C=3C\) دارد. در افت
خروجی، مسیر تخلیه از \lr{nMOS} می‌گذرد و در صعود خروجی، مسیر
شارژ از \lr{pMOS} می‌گذرد. مقاومت هر دو مسیر در این مدل \(R\)
است، زیرا مقاومت \lr{pMOS} با عرض ۲ برابر \(2R/2\) می‌شود.
دو ناحیه \lr{Drain} متصل به خروجی در تقریب ساده ظرفیت \(3C\)
دارند. اگر خروجی، ورودی یک \lr{Inverter} هم‌اندازه را نیز
راه‌اندازی کند و ظرفیت سیم قابل صرف‌نظر باشد، بار گره \(6C\)
است \cite[شکل \secnum{۴.۶}، صفحه ۱۴۷]{weste}. در گیت‌هایی
با ترانزیستورهای سری، مقاومت مسیر از جمع مقاومت ترانزیستورهای
روشن به دست می‌آید. گره‌های میان آن‌ها نیز می‌توانند ظرفیت مؤثر
در گذار داشته باشند.

\subsection{پاسخ تک‌خازنی و قرارداد زمان نصف تغذیه}
\label{sec:delay-rc-half-time}

ابتدا یک مقاومت خطی \(R_{\mathrm{phys}}\) و یک خازن \(C_\Sigma\)
را در نظر بگیرید. با اعمال پله ایدئال در \(t=0\)،
\NotesTopic{eq:foundation-rc}{پاسخ \lr{RC} فصل پیش‌نیاز}
برای شارژ از صفر به \(V_{DD}\) قابل استفاده است. اگر خازن ابتدا
در \(V_{DD}\) باشد و از راه مقاومت به زمین تخلیه شود، در معادله
گره به‌جای \(V_{DD}-v(t)\)، اختلاف ولتاژ تخلیه \(v(t)\) قرار
می‌گیرد و جریان از خازن خارج می‌شود. با تعریف ثابت زمانی
\(\tau_{\mathrm{RC}}=R_{\mathrm{phys}}C_\Sigma\)، دو پاسخ چنین‌اند
\cite[بخش \secnum{۴.۳.۴}، صفحات ۱۴۸ و ۱۴۹]{weste}.
\begin{equation}
  \begin{aligned}
    v_{\mathrm{rise}}(t)&=V_{DD}\bigl(1-e^{-t/\tau_{\mathrm{RC}}}\bigr),\\
    v_{\mathrm{fall}}(t)&=V_{DD}e^{-t/\tau_{\mathrm{RC}}},
      \qquad t\ge 0.
  \end{aligned}
  \label{eq:delay-rc-step}
\end{equation}
در \(t=0\)، ولتاژ شارژ صفر و ولتاژ تخلیه \(V_{DD}\) است.
در هر دو پاسخ، شرط \(v(t_{50})=V_{DD}/2\) نتیجه می‌دهد که
\(e^{-t_{50}/\tau_{\mathrm{RC}}}=1/2\). با گرفتن لگاریتم طبیعی
از دو طرف و استفاده از \(\ln(1/2)=-\ln 2\) داریم
\begin{equation}
  t_{50}=R_{\mathrm{phys}}C_\Sigma\ln 2.
  \label{eq:delay-rc-physical-half-time}
\end{equation}
برای نمونه، \(R_{\mathrm{phys}}=\qty{2}{\kilo\ohm}\) و
\(C_\Sigma=\qty{10}{\femto\farad}\) ثابت زمانی
\qty{20}{\pico\second} و \(t_{50}\approx\qty{13.9}{\pico\second}\)
می‌دهند. این اعداد فرضی‌اند. اگر بار دو برابر شود و مقاومت تغییر
نکند، زمان رسیدن به نصف تغذیه نیز دو برابر می‌شود.

کتاب ضریب \(\ln 2\) را در مقاومت مؤثرِ مدل زمان جذب می‌کند.
اگر \(R_{\mathrm{eff}}=R_{\mathrm{phys}}\ln 2\) انتخاب شود،
رابطه \eqref{eq:delay-rc-physical-half-time} به
\(t_{50}=R_{\mathrm{eff}}C_\Sigma\) تبدیل می‌شود. کتاب سپس
برای ساده‌نویسی \(R_{\mathrm{eff}}\) را با \(R\) می‌نویسد.
در این قرارداد، برای \lr{Inverter} مرجعِ بخش
\ref{sec:delay-rc-sizing} با بار کل \(6C\)، تخمین تأخیر
\(6RC\) است و ضریب \(\ln 2\) پیش‌تر در \(R\) منظور شده
است \cite[رابطه \secnum{۴.۹} و شکل \secnum{۴.۶}، صفحات
۱۴۷ و ۱۴۹]{weste}.

\subsection{از یک گره تا درخت \lr{RC}}
\label{sec:delay-elmore}

وقتی سیم یا شاخه‌های گیت مقاومت قابل‌توجه دارند، گره‌ها می‌توانند با
سرعت‌های متفاوت شارژ شوند. شکل \ref{fig:delay-elmore-tree} نمونه‌ای
از درخت \lr{RC} با یک مسیر مشترک و دو شاخه است. درخت حلقه ندارد،
پس از منبع تا هر گره تنها یک مسیر وجود دارد. در شکل، \(Y\) خروجی
موردنظر و \(Z\) انتهای شاخه دیگر است. \(C_X\)، \(C_Y\) و
\(C_Z\) به‌ترتیب ظرفیت‌های گره‌های \(X\)، \(Y\) و \(Z\) هستند.

\begin{figure}[htbp]
  \centering
  \begin{latin}
  \begin{circuitikz}
    \draw (0,0) node[ocirc,label=above:\(V_s(t)\)] {}
      to[R,l=\(R_0\)] (3,0) node[circ,label=above left:\(X\)] {}
      to[R,l=\(R_1\)] (6,0) node[circ,label=right:\(Y\)] {}
      to[C,l=\(C_Y\)] (6,-1.8) node[ground] {};
    \draw (3,0) to[C,l=\(C_X\)] (3,-1.8) node[ground] {};
    \draw (3,0) -- (3,1.45)
      to[R,l=\(R_2\)] (8.5,1.45) node[circ,label=right:\(Z\)] {}
      to[C,l=\(C_Z\)] (8.5,-1.8) node[ground] {};
  \end{circuitikz}
  \end{latin}
  \caption{درخت \lr{RC} برای محاسبه تأخیر از منبع \(V_s(t)\) تا \(Y\).
  بار شاخه \(Z\) نیز از مقاومت مشترک \(R_0\) شارژ می‌شود.
  شکل برای مثال آموزشی و بر اساس روش بخش \secnum{۴.۳.۵} کتاب
  رسم شده است \cite{weste}.}
  \label{fig:delay-elmore-tree}
\end{figure}

برای گره هدف \(q\)، از بین مقاومت‌های مسیر منبع تا گره \(i\)،
فقط مقاومت‌هایی را جمع می‌کنیم که در مسیر منبع تا \(q\) نیز
هستند. این جمع \(R_{\mathrm{common}}(i,q)\) است. اگر \(C_i\)
ظرفیت گره \(i\) باشد، ثابت زمانی \lr{Elmore} برای گره هدف از
جمع وزن‌دار زیر به دست می‌آید \cite[روابط \secnum{۴.۱۳}
و \secnum{۴.۱۴}، صفحه ۱۵۰]{weste}.
\begin{equation}
  \tau_{\mathrm{E},q}=\sum_i R_{\mathrm{common}}(i,q)C_i.
  \label{eq:delay-elmore-common-path}
\end{equation}
برای فهم وزن‌ها، فرض کنید منبع از صفر به یک ولتاژ ثابت پله می‌زند
و همه خازن‌ها ابتدا تخلیه‌اند. بار نهایی هر خازن برابر ظرفیتِ آن
ضرب‌در تغییر ولتاژ منبع است. تمام این بار از مقاومت‌های میان
منبع و آن خازن عبور می‌کند. از سوی دیگر، اختلاف ولتاژ لحظه‌ای
بین منبع و گره هدف فقط از افت مقاومت‌های مسیر همان گره تشکیل
می‌شود. بنابراین برای محاسبه مساحت زیر منحنی اختلاف ولتاژ منبع
و گره هدف، سهم هر بار تنها از مقاومت‌های مشترک وارد می‌شود.
تقسیم این مساحت بر دامنه پله منبع، رابطه
\eqref{eq:delay-elmore-common-path} را با واحد زمان می‌دهد.

برای خروجی \(Y\) در شکل \ref{fig:delay-elmore-tree}، مسیر منبع
تا \(Y\) شامل \(R_0\) و \(R_1\) است. بارهای \(X\) و \(Z\)
از \(R_0\) سهم می‌گیرند و بار \(Y\) از هر دو مقاومت. در نتیجه
\begin{equation}
  \tau_{\mathrm{E},Y}
    =R_0C_X+(R_0+R_1)C_Y+R_0C_Z
    =R_0(C_X+C_Y+C_Z)+R_1C_Y.
  \label{eq:delay-elmore-branch-example}
\end{equation}
مقاومت \(R_2\) در مسیر منبع تا \(Y\) قرار ندارد. مقدار آن
می‌تواند شکل موج واقعی را تغییر دهد، ولی در این جمع ظاهر
نمی‌شود. با حذف شاخه \(Z\)، رابطه به فرم مدار دوگرهی کتاب
کاهش می‌یابد \cite[شکل \secnum{۴.۱۰} و مثال
\secnum{۴.۳}، صفحات ۱۴۹ و ۱۵۰]{weste}.

در یک مثال فرضی، \(R_0=\qty{1}{\kilo\ohm}\)،
\(R_1=\qty{2}{\kilo\ohm}\)، \(R_2=\qty{3}{\kilo\ohm}\) و
\(C_X\)، \(C_Y\) و \(C_Z\) به‌ترتیب \qty{10}{\femto\farad}،
\qty{20}{\femto\farad} و \qty{30}{\femto\farad} هستند. با
توجه به اینکه حاصل‌ضرب یک \unit{\kilo\ohm} در یک
\unit{\femto\farad} برابر یک \unit{\pico\second} است، رابطه
\eqref{eq:delay-elmore-branch-example} می‌دهد
\[
  \tau_{\mathrm{E},Y}
    =[1(10+20+30)+2(20)]\,\unit{\pico\second}
    =\qty{100}{\pico\second}.
\]
اگر پاسخ را با یک منحنی نمایی با همین ثابت زمانی تقریب بزنیم،
\(t_{50}\approx\tau_{\mathrm{E},Y}\ln 2\) و در این مثال
\(t_{50}\approx\qty{69.3}{\pico\second}\) است. این عدد
از تقریب تک‌نمایی به دست آمده و شکل موج درخت به‌طور کامل حل
نشده است. کتاب در محاسبه‌های دستی، ضریب \(\ln 2\) را در
مقاومت‌ها منظور می‌کند و همان جمع وزن‌دار را برای تخمین تأخیر
به کار می‌برد \cite[بخش‌های \secnum{۴.۳.۴} و
\secnum{۴.۳.۵}، صفحات ۱۴۹ و ۱۵۰]{weste}.

\subsection{اثر \lr{Layout} و تعیین مقاومت مدل}
\label{sec:delay-rc-calibration}

فرض تناسب ظرفیت \lr{Diffusion} با عرض، برای تخمین اولیه مفید
است. پس از تعیین \lr{Layout}، باید ناحیه‌های مشترک و اتصال‌های
واقعی را شمرد. در نمونه \lr{NAND} سه‌ورودی کتاب، دو \lr{Drain}
در شبکه \lr{pMOS} یک ناحیه تماس‌دار مشترک دارند. ظرفیت تخمینی
خروجی در آن طرح از \(9C\) به \(7C\) می‌رسد. تقسیم یک ترانزیستور
پهن به چند شاخه موازی نیز، بسته به چیدمان و اشتراک ناحیه‌ها،
می‌تواند ظرفیت را تغییر دهد \cite[بخش \secnum{۴.۳.۶} و
شکل‌های \secnum{۴.۱۷} و \secnum{۴.۱۸}، صفحات ۱۵۳ و ۱۵۴]{weste}.

برای تخمین مقاومت مؤثر از جریان، تخلیه خروجی از \(V_{DD}\) تا
\(V_{DD}/2\) را در نظر بگیرید. اگر \(I_{\mathrm{eff}}\) جریان
نماینده در همین بازه باشد، باید بار
\(C_\Sigma V_{DD}/2\) جابه‌جا شود. تقسیم بار بر جریان، زمان
تخمینی را می‌دهد. یکسان‌گرفتن این زمان با
\(R_{\mathrm{eff}}C_\Sigma\) نتیجه می‌دهد
\begin{equation}
  \label{eq:delay-rc-current-calibration}
  t_{50}\approx\frac{C_\Sigma V_{DD}}{2I_{\mathrm{eff}}},
  \qquad R_{\mathrm{eff}}\approx\frac{V_{DD}}{2I_{\mathrm{eff}}}.
\end{equation}
در تخمین \lr{Ramp} کتاب، جریان نماینده را از میانگین جریان
ابتدا و انتهای بازه می‌گیرند. شرایط بایاس، شیب ورودی و جهت
گذار در انتخاب مقدار مؤثر نقش دارند. مقدار دقیق‌تر را می‌توان
با مدل افزاره و شبیه‌سازیِ همان گذار به دست آورد
\cite[بخش \secnum{۴.۳.۷} و روابط \secnum{۴.۱۷} تا
\secnum{۴.۱۹}، صفحات ۱۵۴ و ۱۵۵]{weste}.

مدل \lr{RC} نشان می‌دهد چرا بار بیشتر، مسیر مقاوم‌تر و ظرفیت
شاخه‌های مشترک می‌توانند تأخیر را افزایش دهند. در مدل خطی تأخیر،
همین رفتار را برای مقایسه گیت‌ها با دو سهمِ وابسته به بار و
وابسته به ظرفیت‌های درونی بیان می‌کنیم.
